% ECONOMICS_PROBLEM: {"id": "C79-657", "title": "Encoding and complexity of token-rich Battle Sheep", "jel_code": "C79", "source_id": "OP-0221"}
% !TeX program = xelatex
\documentclass[11pt,a4paper]{article}
\usepackage[margin=23mm]{geometry}
\usepackage{fontspec}
\setmainfont{Latin Modern Roman}
\usepackage{amsmath,amssymb,mathtools}
\usepackage{xcolor,enumitem,booktabs,longtable,array,xurl,hyperref}
\definecolor{muted}{HTML}{53636C}
\hypersetup{colorlinks=true,urlcolor=blue,linkcolor=blue}
\setlength{\parindent}{0pt}
\setlength{\parskip}{5pt}
\newcommand{\R}{\mathbb R}
\newcommand{\E}{\mathbb E}
\newcommand{\N}{\mathbb N}
\newcommand{\ind}{\mathbf 1}
\newcommand{\Var}{\operatorname{Var}}
\newcommand{\Cov}{\operatorname{Cov}}
\newcommand{\supp}{\operatorname{supp}}
\DeclareMathOperator*{\argmin}{arg\,min}
\DeclareMathOperator*{\argmax}{arg\,max}
\newcommand{\entrymeta}[3]{{\sffamily\small\color{muted}\textbf{\texttt{#1}}\quad|\quad #2\par #3\par}}
\newcommand{\Problemref}[1]{\texttt{#1}}
\newcommand{\JELref}[2]{\texttt{#2} (JEL \texttt{#1})}
\begin{document}
\section*{Encoding and complexity of token-rich Battle Sheep}
\textbf{Problem ID:} \texttt{C79-657}\quad \textbf{JEL:} \texttt{C79}\par
% BEGIN ECONOMICS STATEMENT
\label{problem:OP-0221}
% GAME_THEORY_PROBLEM: {"local_id": "GT-091", "original_number": 91, "kind": "R", "entry_kind": "statement_only", "published": false, "review_required": true, "category": "game_theory"}
\label{game:GT-091}
{\sffamily\small\color{muted}
\noindent\textbf{ID:} GT-091. \textbf{Category:} Complexity of combinatorial games.
\textbf{Status:} R; source-stated question retained with a rule and encoding audit, not certified as an unresolved classification.
\par}
\subsection*{Definitions and mathematical statement}
Use the finite-board, two-player, normal-play version of generalized Battle Sheep. The input explicitly lists a finite set $B$ of hexagonal-grid cells, stack owners, positive integer stack sizes written in binary, and the player to move. Missing cells are board boundaries. A legal move chooses a stack of the moving player's sheep of size $s\geq2$, chooses an integer $q$ with $1\leq q<s$, and chooses one of the six straight grid directions having an adjacent empty board cell. The $q$ sheep travel along that direction through consecutive empty cells to the last empty cell before a boundary or occupied cell. They form a new stack there, and $s-q\geq1$ sheep remain at the source. Stacks cannot pass occupied cells. A player without a legal move loses.

Let $\mathcal W_{\mathrm{BS}}$ be the language of such explicitly encoded positions from which the player to move can force a win. The catalog question asks how its complexity changes when stack sizes, and hence the total number of sheep, can be exponential in $|B|$.

An essential scope observation is that every legal move occupies one previously empty cell and empties no cell. If $e_0$ cells are initially empty, every play has at most $e_0$ moves, independently of the number of sheep. Moreover every stack size has at most the bit length of the conserved total. Thus the explicit-board, binary-stack model admits a polynomial-space game-tree evaluation. The correct classification target must take this observation into account; merely increasing stack sizes does not supply exponential play depth.

\subsection*{Short English statement}
Check whether large binary-encoded sheep stacks actually change winner-determination complexity once the splitting and movement rules are enforced.

\subsection*{Variants and scope boundaries}
The source's Open Problem 1 asks about EXPTIME-hardness with exponential stack sizes. In the explicit-board model above, EXPTIME-hardness under polynomial-time many-one reductions would imply $\mathsf{PSPACE}=\mathsf{EXPTIME}$. A succinctly encoded exponentially large board, moves that empty the source, or moves onto occupied cells would require a different analysis. These variants are not inferred from the original rules. The polynomial depth observation does not bound the number of legal splits by a polynomial.

\subsection*{Source relationship and status}
The source proves PSPACE-completeness for its restricted family and explicitly states Open Problem 1 for exponentially large stacks. The occupancy observation is an editorial audit of Section 1.1's rules, not a claimed published resolution. Accordingly this entry preserves the historical question while declining to repeat the catalog's claim that token magnitude defeats polynomial-depth membership. Combining the audit with the source hardness result gives PSPACE-completeness for the explicit-board extension, provided the restricted instances use the same winner convention.

\subsection*{References}
\begingroup\footnotesize\raggedright
\noindent [S1] Kyle Burke and Hirotaka Ono, \emph{Battle Sheep is PSPACE-complete} (2025), arXiv:2505.06414v1. Inspected locators: Section 1.1 (bottom sheep stays, empty-cell movement), Theorem 1 and Corollary 1, Section 3, Open Problem 1.
\url{https://arxiv.org/html/2505.06414v1}
\par\endgroup

\subsection*{Common conventions: Source mathematical conventions}

\textbf{Finite games.} For a finite set $A$, $\Delta(A)$ is its probability
simplex. Players choose independent mixed strategies in Nash equilibrium;
correlation is allowed only when explicitly part of the solution concept.
In a game with payoff functions $u_i$ and strategy profile $x$, an additive
$\varepsilon$ Nash equilibrium satisfies
\[
 u_i(x)\geq u_i(x_i',x_{-i})-\varepsilon
 \quad\text{for every player $i$ and every admissible deviation $x_i'$}.
\]
For numerical approximation, payoffs are normalized as locally specified.
An $\varepsilon$ payoff residual is distinct from distance at most
$\varepsilon$ to the set of exact equilibria. Point convergence, convergence
of empirical play, and convergence in distance to an equilibrium set are
also distinct.

\textbf{Computational representations.} Unless stated otherwise, a complexity
question takes rational data with binary encoding; polynomial time is measured
in total bit length. A fixed-accuracy polynomial algorithm is not necessarily
polynomial in $1/\varepsilon$, and neither is necessarily polynomial in
$\log(1/\varepsilon)$. Succinct games and value, demand, or sampling oracles
must declare their interfaces. Exact real solutions require an explicit
representation; an unspecified list of arbitrary real numbers is not a
finite computational output. A claimed algorithmic barrier is meaningful
only for its specified model and reductions.

\textbf{Dynamic games.} Histories, information, observation, and allowable
randomization are part of the model. A stationary strategy depends only on
the current state; a positional strategy in a deterministic graph game
chooses one action at each controlled vertex. Nash equilibrium at an initial
state and equilibrium after every history are different requirements.
An expectation of a pathwise $\liminf$ payoff, a $\liminf$ of expected averages,
a limiting expected average, and a uniform finite-horizon equilibrium are
different criteria. None is silently substituted for another.

\textbf{Allocation and mechanisms.} A complete allocation assigns every item
exactly once unless otherwise stated. Zero-valued goods, negative values,
capacity constraints, feasibility, lotteries, and copies of goods affect
fairness definitions and are specified locally. Dominant-strategy incentive
compatibility, Bayesian incentive compatibility, universal truthfulness,
and truthfulness in expectation impose different inequalities. Whenever
an objective ratio has zero denominator, its convention is stated or those
instances are excluded.

\textbf{Infinite-dimensional models.} Expectations, integrals, stochastic
equations, and optimization objectives require their local regularity,
integrability, filtrations, and admissible domains. A backward--forward
mean-field system is a boundary-value problem; it is not an initial-value
semiflow without an additional construction. Long-time, large-population,
and vanishing-noise limits require an order of limits and a convergence
topology. If a minimum need not be attained, the objective is its infimum
or supremum and the approximation requirement is stated separately.
% END ECONOMICS STATEMENT
\end{document}
